\chapter{Multi-User Functionality}

The ability to support multiple users within a shared virtual environment is one of the distinguishing features of the proposed virtual laboratory. Unlike traditional virtual laboratory applications designed for individual learning, the implemented system enables several participants to join the same laboratory session, observe each other's activities, and collaborate during experimental tasks. This collaborative capability enhances social interaction, supports instructor-led demonstrations, and creates learning experiences that more closely resemble conventional laboratory classes.

The multi-user functionality is implemented using a client--server architecture in which a centralized server maintains the authoritative state of the virtual environment while synchronizing user actions across all connected participants. Every user is represented by an avatar, allowing participants to observe movements and interactions within the shared laboratory. Experimental objects, user positions, and interaction events are synchronized in real time to ensure that all participants experience a consistent representation of the virtual environment.

The implementation was designed to support collaborative experimentation without affecting the usability of the virtual laboratory. Users can freely navigate between different laboratory areas, participate in shared experiments, and observe changes made by other participants during the same laboratory session. This functionality enables collaborative learning, instructor-guided demonstrations, and group discussions within an immersive Virtual Reality environment.

This chapter describes the implementation of the collaborative features developed for the virtual laboratory, including user connection management, avatar representation, synchronization of user interactions, and the mechanisms that maintain a consistent shared virtual environment. The chapter also discusses the practical challenges encountered during the implementation of the multiplayer system and the limitations of the current approach.

\section{Collaborative Interaction}

Collaborative interaction is a fundamental feature of the proposed virtual laboratory, enabling multiple users to participate simultaneously within the same immersive environment. Unlike single-user virtual laboratories, the implemented system allows participants to observe each other's presence, communicate through shared visual interactions, and collaborate while performing laboratory experiments. This shared experience promotes active learning and closely resembles practical laboratory sessions conducted in educational institutions.

When users join a laboratory session, they are connected to a common virtual room where each participant is represented by an avatar. The avatar reflects the user's position and orientation within the virtual environment, allowing other participants to observe movement and activity in real time. This representation improves spatial awareness and enhances the sense of social presence during collaborative learning activities.

Users can independently navigate throughout the virtual laboratory while remaining connected to the same collaborative session. They are free to move between the classroom, laboratory, corridor, and other learning spaces without interrupting the experience of other participants. Since all users share the same virtual environment, discussions and demonstrations can naturally occur while exploring different parts of the laboratory.

Collaborative interaction extends beyond user movement to include shared manipulation of experimental equipment. When a participant interacts with laboratory objects, such as repositioning the convex lens or adjusting the projection screen, the corresponding changes are synchronized and immediately become visible to all connected users. Consequently, every participant observes the same experimental configuration, supporting collaborative experimentation and instructor-guided demonstrations.

The shared interaction model also enables multiple learners to discuss experimental observations while viewing identical experimental results. Students can compare different optical configurations, observe the effects of parameter changes, and collectively analyse the behaviour of the virtual experiment. This collaborative approach encourages discussion, peer learning, and active participation throughout the laboratory session.

The collaborative interaction mechanisms were implemented using the multiplayer framework described in Chapter~4 and integrated into the browser-based virtual laboratory without altering the interaction model for individual users. As a result, both desktop and Virtual Reality users can participate within the same collaborative environment regardless of their hardware platform. The implementation of the collaborative interaction system is available in the accompanying project repository \cite{thesisRepo}.

\section{Implementation}

The multi-user functionality was implemented using a client--server architecture that enables multiple users to participate simultaneously within the same virtual laboratory. The implementation builds upon the browser-based architecture described in Chapter~4 by integrating real-time networking, avatar synchronization, and shared experiment interaction into a unified collaborative environment. The objective of the implementation was to provide a consistent and responsive shared experience while maintaining compatibility with both desktop browsers and Virtual Reality (VR) devices.

Each user connects to the laboratory through the multiplayer server, which assigns the participant to a collaborative session. Once connected, an avatar representing the user is created within the virtual environment. The avatar continuously updates its position and orientation according to the user's movement, allowing other participants to observe navigation and interaction throughout the laboratory. The synchronization process ensures that every connected client maintains a consistent representation of all participants within the same session.

Shared laboratory objects are synchronized using an authoritative server model. Whenever a user manipulates an experimental component, such as moving the convex lens, adjusting the projection screen, or modifying the focal length, the corresponding state changes are transmitted to the server. After validating the updated state, the server distributes the changes to all connected clients, ensuring that every participant observes an identical experimental configuration. This synchronization mechanism allows collaborative experimentation without inconsistencies between users.

The implementation also supports collaborative observation and instructor-guided demonstrations. Since all participants share the same virtual environment, instructors can demonstrate experiments while students observe the resulting changes in real time. Likewise, students can discuss experimental observations, compare different configurations, and collaboratively investigate physical phenomena without requiring physical access to laboratory equipment.

To maintain responsive interaction, only the information required to synchronize the collaborative environment is exchanged between the client and the server. User transformations, avatar updates, object movements, and interaction events are transmitted as incremental state updates, while static assets such as three-dimensional models, textures, and environment geometry remain locally available on each client. This approach minimizes communication overhead while maintaining smooth interaction during collaborative laboratory sessions.

The modular implementation adopted throughout the virtual laboratory also facilitates future extension of the multiplayer system. Additional collaborative experiments, interaction techniques, user roles, and educational features can be incorporated using the existing networking infrastructure without requiring major architectural modifications. The complete implementation of the multiplayer functionality is available in the accompanying project repository \cite{thesisRepo}.

\section{Challenges and Limitations}

Although the proposed virtual laboratory successfully demonstrates browser-based Virtual Reality (VR), collaborative learning, and interactive physics experiments, several challenges were encountered during its development. These challenges primarily relate to browser-based VR technologies, multiplayer synchronization, hardware compatibility, and the scope of the implemented educational content.

One of the primary challenges was achieving smooth real-time synchronization between multiple users. Since collaborative interaction depends on continuous communication between clients and the server, network latency and connection quality can influence the responsiveness of synchronized user movements and shared experimental objects. Although the implemented synchronization mechanism provides a consistent shared environment under normal network conditions, performance may be affected when users experience unstable or high-latency Internet connections.

Another challenge arises from the diversity of hardware and software platforms supporting WebXR. While the application was primarily developed and tested using the Meta Quest 3 headset and modern desktop browsers, browser implementations and WebXR support may differ across devices. Consequently, user experience and performance may vary depending on the capabilities of the target hardware and browser.

The browser-based architecture also introduces performance limitations compared with native VR applications. Rendering complex three-dimensional scenes while simultaneously processing user interaction and multiplayer synchronization places additional demands on the client device. Although the current implementation performs satisfactorily for the developed virtual laboratory, larger environments or more computationally intensive simulations may require additional optimization techniques such as level-of-detail rendering, model simplification, asset streaming, or more advanced synchronization strategies.

The current implementation focuses on two representative physics experiments: the Stern--Gerlach experiment and the Convex Lens experiment. While these experiments demonstrate the flexibility of the proposed virtual laboratory framework, they represent only a small subset of the experiments typically encountered in undergraduate physics education. The modular architecture adopted in this thesis provides a foundation for integrating additional experiments and educational modules in future work.

Finally, the evaluation presented in this thesis concentrates on the technical implementation and functionality of the virtual laboratory. A comprehensive educational evaluation involving students and instructors was outside the scope of this work. Future research should investigate learning effectiveness, usability, user experience, and collaborative learning outcomes through structured user studies conducted in educational environments.

Despite these limitations, the developed virtual laboratory demonstrates that modern web technologies can be successfully combined with immersive Virtual Reality and real-time collaboration to create an accessible platform for interactive physics education. The modular implementation also provides a solid foundation for future expansion and continued development of browser-based collaborative virtual laboratories.

\section{Summary}

This chapter presented the implementation of the multi-user functionality developed for the browser-based virtual laboratory. The collaborative features enable multiple participants to join the same laboratory session, navigate the virtual environment, observe each other's presence through avatars, and interact with shared experimental equipment in real time. These capabilities extend the virtual laboratory beyond a conventional single-user application and provide an immersive environment that supports collaborative learning.

The implementation integrates user connection management, avatar representation, and real-time synchronization within a unified client--server architecture. By maintaining a consistent shared state across all connected participants, the system enables collaborative experimentation and instructor-guided demonstrations while supporting both desktop and Virtual Reality users through a common browser-based platform.

The chapter also discussed the practical challenges encountered during development, including network latency, browser compatibility, hardware limitations, and the restricted scope of the implemented experiments. Although these limitations provide opportunities for future improvements, the current implementation demonstrates the feasibility of combining modern web technologies, immersive Virtual Reality, and real-time networking to create an accessible collaborative learning environment.

Overall, the multi-user functionality represents one of the key contributions of this thesis, enabling collaborative physics education within a browser-based virtual laboratory. Together with the interactive experiments presented in the previous chapter, the implemented system establishes a foundation for future extensions involving additional educational content, enhanced collaboration features, and broader deployment in academic environments.